\chapter{Типизированные внутренние кандидаты карты c0}
\label{ch:v8-baryon-c0-typed-internal-map-candidate-audit-gate}

Для однополюсного барионного ядра вспомогательное поле должно иметь
квадратичный блок
\begin{equation}
 S_{\rm aux}=\frac12\phi(z+m^2)\phi-g\phi J,
 \qquad m^2=c_0a.
 \label{eq:v8-baryon-c0-required-auxiliary-block}
\end{equation}
Следовательно, кандидат на $c_0$ должен быть не просто положительным
безразмерным числом. Требуются внутренний выбор, независимость от целевой
массы, объявленный морфизм носителей и появление именно в гессиане поля
$\phi$. Обозначим условие допуска
\begin{equation}
 A(x)=D(x)S(x)T(x)H(x)N(x),
 \label{eq:v8-baryon-c0-admission-predicate}
\end{equation}
где $D$ означает безразмерность, $S$ --- внутренний выбор, $T$ ---
типизированную карту в auxiliary carrier, $H$ --- появление в его
квадратичном гессиане, а $N$ --- независимость от наблюдаемой мишени.

\section{Конечный реестр кандидатов}

База формул поставляет семь ближайших объектов:
\begin{equation}
 \mathcal C_{c_0}=
 \left\{1,3,4,\frac38,e^{-2},R_\chi,\mathcal R_3\right\},
 \qquad
 \mathcal R_3:=\frac{\lambda_3^2}{\alpha\beta}.
 \label{eq:v8-baryon-c0-candidate-set}
\end{equation}
Первые три являются finite-Dirac, vector и scalar mass-square ratios
Version III; $3/8$ есть спектральный gauge coupling; $e^{-2}$ выбрано
индексом цепного процесса; $R_\chi=1$ только на дополнительно выбранном
чистом тепловом профиле; $\mathcal R_3=4$ возникает из следовой нормы
сдвинутой кривизны. Их точные значения в соответствующих ветвях равны
\begin{equation}
 (c_f,c_A,c_s,g^2,r_{\rm KMS},R_\chi,\mathcal R_3)
 =\left(1,3,4,\frac38,e^{-2},1,4\right).
 \label{eq:v8-baryon-c0-candidate-values}
\end{equation}

Первые шесть объектов принадлежат hidden finite, gauge, Markov или
spectral-profile носителям. В проекте нет объявленной стрелки
\begin{equation}
 \mathcal M_x:\mathcal C_x
 \dashrightarrow
 \operatorname{End}(\mathcal H_{\rm aux}),
 \qquad
 \mathcal M_x(x)=c_0I_{\rm aux},
 \label{eq:v8-baryon-c0-missing-carrier-map}
\end{equation}
которая помещала бы любой из них в коэффициент $m^2=c_0a$.

\section{Сильнейший близкий кандидат}

Сдвинутая кривизна находится на том же следовом носителе, который породил
$W_3$, и потому является единственным близким кандидатом. Но она даёт
\begin{equation}
 \mathcal R_3=\frac{\lambda_3^2}{\alpha\beta}=4,
 \qquad
 c_0=\frac{m_{\rm pole}^2}{a},
 \label{eq:v8-baryon-c0-shape-versus-pole-ratio}
\end{equation}
то есть отношение трёх коэффициентов взаимодействия, а не отношение
квадратичной массы нового auxiliary carrier к масштабу базы. Центральный
фон, из которого получено первое равенство, вдобавок не является нечётным
оператором Дирака при ненулевом масштабе.

Точный булев реестр имеет вид
\begin{equation}
 (A(c_f),A(c_A),A(c_s),A(g^2),A(r_{\rm KMS}),A(R_\chi),A(\mathcal R_3))
 =(0,0,0,0,0,0,0),
 \label{eq:v8-baryon-c0-admission-ledger}
\end{equation}
и потому
\begin{equation}
 \sum_{x\in\mathcal C_{c_0}}A(x)=0\quad\hbox{из }7.
 \label{eq:v8-baryon-c0-admission-count}
\end{equation}

\begin{theorem}[Исчерпание объявленных внутренних отношений]
Ни одно из семи уже вычисленных или условно выбранных безразмерных отношений
не является типизированной картой $c_0$ в гессиан вспомогательного поля
барионного ядра. Результат является исчерпанием текущего именованного
реестра, а не запретом на существование нового cross-carrier морфизма.
\end{theorem}

\begin{nogocriterion}[Численное совпадение не является морфизмом]
Совпадение скалярного mass-square ratio и $\mathcal R_3$ со значением $4$
не разрешает положить $c_0=4$. Их источники, роли в действии и носители
различны; требуется явная вариация одного родителя, дающая коэффициент
$c_0a$ в auxiliary Hessian.
\end{nogocriterion}

\begin{protocol}\begin{enumerate}
\item проверено кандидатов: \textbf{$7$};
\item прошедших полный контракт: \textbf{$0$};
\item сильнейший близкий кандидат: \textbf{$\mathcal R_3=4$};
\item дефект сильнейшего кандидата: \textbf{shape ratio, не pole ratio};
\item универсальный запрет будущих карт: \textbf{не утверждается};
\item следующий гейт: \textbf{минимальная архитектура cross-carrier морфизма}.
\end{enumerate}\end{protocol}

Машинный сертификат сохранён в
\path{s2t_v8_baryon_c0_typed_internal_map_candidate_audit_gate_results.json}.